University of Wollongong
Research Online
Faculty of Informatics - Papers (Archive) Faculty of Engineering and Information
Sciences
1-1-2007
Combining i* and BPMN for business process model lifecycle management
George Koliadis
University of Wollongong, gk56@uowmail.edu.au
Aleksander Vranesevic
University of Wollongong, av85@uow.edu.au
Moshiur Bhuiyan
University of Wollongong, mmrb95@uow.edu.au
Aneesh Krishna
University of Wollongong, aneesh@uow.edu.au
Aditya K. Ghose
University of Wollongong, aditya@uow.edu.au
Follow this and additional works at: https://ro.uow.edu.au/infopapers
Part of the Physical Sciences and Mathematics Commons
Recommended Citation
Koliadis, George; Vranesevic, Aleksander; Bhuiyan, Moshiur; Krishna, Aneesh; and Ghose, Aditya K.:
Combining i* and BPMN for business process model lifecycle management 2007.
https://ro.uow.edu.au/infopapers/575
Research Online is the open access institutional repository for the University of Wollongong. For further information
contact the UOW Library: research-pubs@uow.edu.au 
Combining i* and BPMN for business process model lifecycle management
Abstract
The premise behind ‘third wave’ Business Process Management (BPM1) is effective support for change at
levels. Business Process Modeling (BPM2) notations such as BPMN are used to effectively conceptualize
and communicate process configurations to relevant stakeholders. In this paper we argue that the
management of change throughout the business process model lifecycle requires greater conceptual
support achieved via a combination of complementary notations. As such the focus in this paper is on the
co-evolution of operational (BPMN) and organizational (i*) models. Our intent is to provide a way of
expressing changes, which arise in one model, effectively in the other model. We present constrained
development methodologies capable of guiding an analyst when reflecting changes from an i* model to a
BPMN model and vice-versa. 1 Introduction
Disciplines
Physical Sciences and Mathematics
Publication Details
This conference paper was originally published as Ghose, AK, Koliadis, G, Vranesevic, A, Bhuiyan, M and
Krishna, A, Combining i* and BPMN for business process model lifecycle management, in Proceedings of
the BPM-2006 Workshop on Grid and Peer-to-Peer based Workflows, Lecture Notes in Computing
Science, 4103, pp 416-427, 2006.
This conference paper is available at Research Online: https://ro.uow.edu.au/infopapers/575 
Combining i* and BPMN for
Business Process Model Lifecycle Management
George Koliadis
1
, Aleksandar Vranesevic
1
, Moshiur Bhuiyan
1
, Aneesh Krishna
1
,
and Aditya Ghose
1
1
School of Information Technology and Computer Science (SITACS),
University of Wollongong (UOW), Gwynneville NSW 2522, Australia
{gk56, av85, mmrb95, aneesh, aditya}@uow.edu.au
Abstract. The premise behind ‘third wave’ Business Process Management
(BPM
1
) is effective support for change at levels. Business Process Modeling
(BPM
2
) notations such as BPMN are used to effectively conceptualize and
communicate process configurations to relevant stakeholders. In this paper we
argue that the management of change throughout the business process model
lifecycle requires greater conceptual support achieved via a combination of
complementary notations. As such the focus in this paper is on the co-evolution
of operational (BPMN) and organizational (i*) models. Our intent is to provide
a way of expressing changes, which arise in one model, effectively in the other
model. We present constrained development methodologies capable of guiding
an analyst when reflecting changes from an i* model to a BPMN model and
vice-versa.
1 Introduction
Business process models play a key role in both organizational management [1] [2]
and enterprise information systems development [3]. Many notations developed for
the task of modeling business processes, have their own focus of application and
appropriate audience [4] [5] [6] [7] [8]. High-level conceptual models provide an
understanding of an organization from an intentional and social perspective [9] for
reasoning support during redesign [9]. In comparison, lower-level technical models
are especially suited for applications in the description, execution and simulation of
business processes [8].
Business process development should be based on principled high-level models of
the enterprise and the business context. Commonly, processes are formulated in an
ad-hoc fashion without reference to these high-level models. Some of the most
prominent modeling notations enlisted are focused towards technically-oriented data,
and process modeling notations such as ER, Data-Flow, Systems Flowcharting and
UML and workflow modeling [10]. In this work, we offer constrained development
methodologies to guide development of process models from higher-level conceptual
models. This supports life-cycle management in the following sense: when changes
occur to the high-level model, these can be reflected in the process model, and vice-
versa. In this paper, Section 2 provides a background to business process modeling
with an overview of our chosen notations. Section 3 illustrates concepts/methods
provided in our methodologies (with examples). The paper is concluded in Section 4.
2 Background
The notations used for modeling business processes have been categorized in many
works, based on their conceptual features [4] [5] [6] [7] [8]. The common principle
recognized in all analyses is that some notations are more suited towards specific
audiences (i.e. with either technical/non-technical backgrounds) or applications (i.e.
possibly for description, re-design or execution) throughout the business process
lifecycle. Many notations focus on specific aspects, with limited relation/traceability
to other important business process aspects. This has brought about the need for an
enterprise view [6] to support the development and maintenance of rich models that
provide an enhanced ability to conceptualize, communicate and understand business
processes, and their context of operation.
In related work, some preliminary ideas in [11] have been proposed for developing
a BPMN model given the existence, and agreement to, an i* model of the process.
Six steps are provided for mapping between constructs, with no consideration for
reflecting change and consistency made. Also, an approach for deriving a BPMN
model from a business model is proposed in [12], achieved through the intermediate
translation of the business model into an activity dependency model that can then be
translated into a business process model. In this work, we provide a simpler approach
aimed at reducing added complexity and/or misinterpretations during modeling. Furthermore, much work has been completed on supporting guided translation and coevolution of i* into various other behavioral modeling notations and languages [13]
[14] [15]. The primary aim in these approaches is to further develop detailed design
artifacts that can lead onto implemented systems, or directly be used in the configuration of agent-based systems. However, our primary focus is on modeling lifecycle
support during BPM
1
projects whereby the concern is for the development and/or
assessment of detailed business process designs. The work in this paper extends
previous work in [16]. In comparison to previous work, we take the following approach to lifecycle management: when changes to a business process model (i.e.
BPMN – [17]) occur, these changes must ensure some notion of consistency with a
higher-level enterprise model, and vice versa. In this instance, an i* model [9].
2.1 Agent-Oriented Conceptual Modeling (AOCM) with i*
i* supports modeling rich organizational contexts by offering high-level social/anthropomorphic abstractions (such as goals, tasks, soft goals and dependencies)
as modeling constructs for reasoning support during business process redesign [9]
[7]. Figure 1 represents a simple i* Meeting Scheduling model. The central concept
in i* is that of intentional actor. These can be seen in the Meeting Scheduling model
as nodes representing the intentional/social relationships between three (3) actors
required to schedule a meeting: a Meeting Initiator (MI); Meeting Scheduler (MS);
and, Meeting Participant (MP).
Fig. 1. An i* Strategic Rationale (SR) Meeting Scheduling Model with a Routine Illustrated
The i* framework consists of two modeling components [9] Strategic Dependency
(SD), and Strategic Rationale (SR) models. The SD model consists of a set of nodes
and links. Each node represents an actor, and each link between the two actors indicates that one actor depends on the other for something in order that the former may
attain some goal. The depending actor is known as depender, while the actor depended upon is known as the dependee. Dependancies may involve goals to be
achieved (e.g. MeetingBeScheduled), tasks to be performed (e.g. EnterAvailDates),
resources to be furnished (e.g. Agreement), or soft-goals (optimization objectives or
preferences) to be satisficed (e.g. MaximizeAttendance).
The SR mode further represents internal motivations and capabilities (i.e. processes or routines) accessible to specific actors that provide illustration of how dependencies can be met. In i*, a routine [9] specifies an intended course of action an
actor may pursue given a set of alternatives. These elements and their relationships
represent the strategic requirements of a business process when invoked in a specific
context. For example, to ScheduleMeeting (illustrated in Figure 1 with its Scope) that
includes three sub-tasks and six dependencies with two additional actors. Tasks in i*
may be primitively workable whereby the actor responsible for the element believes
that it can achieve its requirements at execution time – i.e. it is sufficiently reduced
during decomposition. In comparison to BPMN however, a primitively workable
element may still be represented as a sub-process as the term does not imply a ‘primitively executable action’ (i.e. application of analyst / designer discretion). Furthermore, for a routine to be workable, all involved actors must be committed to satisfying their dependencies [9].
The Tropos project [18] aims to provide methodological support for advancing the
i* framework further towards architectural and detailed design where dynamic / beRoutine
+ Scope
havioral aspects are of importance. Specifically, Formal Tropos (FT) – see [19], is a
part of the Tropos project that provides a specification language for modeling dynamic aspects of an i* model via formal annotation of Creation and Fulfillment conditions. These conditions are specified using first-order typed linear temporal logic
and prescribe the constraints on an elements lifecycle. In this work, we take the same
approach to annotation (with the use of fulfillment conditions annotated to i* models). In comparison, our work is illustrated via informal annotations.
2.2 Business Process Modeling with BPMN
The Business Process Modeling Notation (BPMN), developed by the Business Process Management Initiative (BPMI.org) [17] is primarily a technically-oriented business process modeling notation that supports the assignment of activity execution
control to entities within an organization via ‘swim-lanes’. BPMN has the capability
to map directly to executable process languages including XPDL [20] and BPEL [17]
[21]. Furthermore, an analysis of BPMN [22] also stated its high maturity in representing concepts required for modeling business process, apart from some limitations
in terms of representing state, and the possible ambiguity of the swim-lane concept.
Fig. 2. A BPMN Patient Treatment Business Process Model
Figure 2 represents a simple BPMN Patient Treatment process. Processes are represented in BPMN using flow nodes: events (circles), activities (rounded boxes), and
decisions (diamonds); connecting objects: control flow links (unbroken directed
lines), and message flow links (broken directed lines); and swim-lanes: pools (highlevel rectangular container), and lanes partitioning pools. These concepts are further
discussed in [17].
3 Constrained Development Methodologies
We propose constrained development methodologies to guide the derivation or maintenance of one type of model given the availability of the other. The development is
supported with the introduction of two concepts: fulfillment conditions (i.e. as in [19])
and effect annotations.
An effect is broadly defined as the result (i.e. product or outcome) of an activity
being executed by some cause or agent. An effect annotation is a specific statement
relating to the outcome of an activity, associated to a state altering construct in a
given model. During BPM
2
, effects are annotated to atomic tasks/activities or subprocesses within an actor’s lane. The execution of a number of activities in succession results in a cumulative effect that includes the specific effects of each activity in
the sequence. We also note the fact that certain effects can undo prior effects (i.e. in
the case of compensatory activities). Effect annotations may possibly be formalized
using the formal layers of some currently well-developed Goal-Oriented Requirements Engineering (GORE) methodologies [23] [19], however, we only state their
applicability in this work, and aim towards possible integration in the future.
Fulfillment conditions are annotated to tasks and goals assigned to actors in an SR
diagram, and dependencies (i.e. not including soft-goals as these are used during
assessment of alternatives and describe non-functional properties to be addressed) in
an i* model. A fulfillment condition [19] is a statement specifying the required conditions realized upon completion of a given task, goal or dependency. Fulfillment conditions recognize the required effects on a business process model. For example, a
fulfillment condition for a task dependency to EnterADateRange, may be the DateRangeCommunicated effect (subsequently required by the task assigned to a dependee
actor).
3.1 Annotation and Propagation
Tasks, goals and dependencies are annotated with fulfillment conditions in an i*
model. Additionally, the tasks assigned to participants in a BPMN model are annotated with effects for assessment against fulfillment conditions.
Tasks associated to dependencies on the dependee side may require additional effects when related to a BPMN model. That is, the fulfillment conditions for a dependency may not be explicitly stated against the tasks. For example, the fulfillment
condition for ProposedDateProvided (i.e. annotated to the ProposedDate resource
dependency in Figure 1) will be propagated to the ObtainAvailDate task. This should
occur during annotation, whenever a fulfillment condition is annotated to a resource,
goal or task dependency.
Effect annotations in BPMN models are propagated via trajectories. A trajectory
is a sequential execution of activities terminating at an end state that represents the
operational goal of the process. Control flow links between events, activities, and
gateways within a BPMN model indicate the flow of trajectories. Effects within a
process are accumulated during forward traversal through a trajectory. This accumu-
lation ensures that any compensatory activities, that may undo effects, are also taken
into account during traversal.
3.1.1 Annotating the Meeting Scheduling Model (Figures 1 and 4).
Table 1. Annotation of Fullfillmnent Conditions to Respective Tasks/Dependancies
Task/Dependency (Figure 1) Fulfillment Conditions Task Annotation (Post
Development – Figure 4)
MI: SchedulerSchedules Meeting DateRangeEnteredIntoScheduler;
DateRangeCommunicatedToScheduler
1;
1;
MS: ScheduleMeeting AgreedDateKnownToInitiator 4
MS: ObtainAvailableDates ProposedDateProvided;
AvailableDatesObtained;
AvailableDatesStored;
AvailableDatesValidated
2 (message);
2;
2;
2
MS: ObtainAgreement AgreementObtained;
AgreementRecorded
4;
4
MS: MergeAvailableDates AvailableDatesMerged 3
P: AgreeToDate DateAgreedTo; AgreementProvided; 6; 6 (message)
P: FindAgreeableDateUsing Scheduler
AvalDatesEnteredIntoScheduler;
AgreeableDateFoundUsingScheduler
5;
6
MS-Dep->MI: EnterDateRange DateRangeCommunicatedToScheduler 1
MI-Dep->MS: MeetingBeScheduled AgreedDateKnownToInitiator 4
MS-Dep->P: EnterAvailDates AvailDatesEnteredIntoScheduler 5
P-Dep->MS: ProposedDate ProposedDateProvided 2
MS-Dep->P: Agreement AgreementProvided 6 (message)
3.2 Scope Projection
In order to evaluate consistency between the two notations, we provide some rules for
projecting the scope of the i* model. In the current case, i* models are likely to represent a broader scope in comparison to a specific BPMN model as they are applied to
capture the greater organizational context. Scope projection is based on an identification of the business process (represented in BPMN) as a routine assigned to an actor
in an i* model.
− Rule 1: The root node of the routine traceable to the process in consideration and
all tasks in its first level of decomposition from are to be within scope.
− Rule 2: All dependencies that are associated to a task within the scope of the routine, where the actor in control of the routine (initiator) is the depender are within
the scope of the process; as well as the tasks assigned to dependee actors.
− Rule 3: All dependencies that are associated to a task within the scope of the routine, where the intiator is the dependee are within the scope of the process iff the
task assigned to the depender is part of some decomposition of a task in the scope
of the process as per Rule 2; as well as the tasks assigned to the depender actors.
3.3 Consistency Evaluation
We introduce consistency rules to provide a mechanism for ensuring consistency
between i* and BPMN models (developed with consideration to [19]).
− Rule 1: Every actor in an i* model required as a participant in the routine (traceable to the business process) and any of their tasks must be represented in the
BPMN model (and vice versa), assessed via application of scope projection rules.
− Rule 2: There must exist a trajectory in the process model, whereby the operational
objective (as encoded in the accumulated fulfillment conditions of traceable tasks)
of the routine is achieved, and the sequence of activities is consistent with the requirements specified in the routine as further outlined below:
− Rule 2.1: The accumulated effect of all tasks and goals traceable to the routine
must achieve accumulated routine fulfillment conditions during forward traversal of at least one trajectory in the process model; AND,
− Rule 2.2: The fulfillment of a task on the depender side of a dependency must
not be realized before the fulfillment of the dependency upon accumulation of
effects during forward traversal of the same trajectory.
3.4 Constrained Development of a Business Process Model given a High-Level
Conceptual Model
These steps are based on the aforementioned consistency rules aimed towards providing analyst guidance during initial model development.
− Step 1: Identify internal and external actors in i* diagram.
− Step 2: Map elements to equivalent constructs within the BPMN model. See substeps below.
− Step 2.1: Map Participants. The greater organization for which the i* model is
represented is signified as a pool in BPMN. Any external participants are also
represented as pools. Internal organizational actors are represented as lanes
within the organizational pool.
− Step 2.2: Map Activities. Tasks within i* are represented as either sub-processes
or atomic activities within BPMN assigned to actors within pools and lanes.
− Step 3: Sequence required tasks/sub-processes and introduce control and sequence
flow links by analyzing fulfillment conditions. Tasks placed within each pool or
lane are now sequenced to conform to routine requirements by taking Consistency
Rule 2 (see: Section 3.3) into consideration. This requires that tasks be sequenced
using control flow links in a manner that results in a trajectory satisfying fulfillment conditions on an i* model. Control flow links are used to indicate realization
of dependencies between actors within the same organization. In order to realize
dependencies between organizational boundaries, a message flow link is used to
represent the dependency going from the depender lane to the dependee lane. This
may require single/multiple messages between tasks derived via analysis of fulfillment conditions.
− Step 4: Elaborate on sub-processes. The choice to introduce tasks or subprocesses into the BPMN diagram for specific tasks in the i* model is made in Step
2.2. The analyst can develop each sub-process guided by the list of required fulfillment conditions annotated to the i* task that the sub-process realizes.
Figure 3 illustrates the application of the constrained development methodology in
the context of the Meeting Scheduling model represented in Figure 1, with annotations applied in Table 1. Much of the detail has been omitted for brevity. The following section describes a possible change requirement and its reflection within an i*
model for further analysis.
Fig. 3. BPMN Process Model derived using the Constrained Development Methodology
3.4.1 Reflecting Changes in an i* Model to an associated BPMN Model.
The scope projection techniques are used to assess whether a change in an i* model
will impact a BPMN model. These guidelines aim to support the reflection of change
between i* and BPMN models for the specific instances of impacting change outlined
below.
− Step 1: For each classification outlined below apply associated changes.
− Addition of an actor. If a new actor has been added to the i* model, a swimlane
(i.e. for an internal actor) or pool (i.e. for an external actor) will need to be
placed on the process model. Additionally, new dependencies must exist between the actor and existing actors (described below). These dependencies will
be included for all new actors where the dependency is related to the routine
and actor is the dependee. However, where the actor is the depender they will
only be included if linked to a task in an existing dependency graph (see Scope
Projection rules).
− Addition of a goal/task/resource dependency. If a new dependency has been
added to the i* model, then this may require the addition of new activities/subprocesses and message flow links within the BPMN model (as described below).
− Addition of a goal or task. The addition of a goal or task will require the addition of a task within the BPMN model. The addition of these tasks must be
Step 1:
Internal /
External
Participant Pools
Step 3:
Task Sequencing,
Message and
Sequence Flow
Step 2:
Participants and
Activities
1
2 3 4
5 6
Step 4:
Sub-Process
Elaboration
scoped to their respective actors, and any dependencies must be realized via
message-flow links where one of the actors is external to the organization.
− Step 2: Re-apply consistency rules to both models to assess whether consistency
has been maintained.
Consider the following example applied to the Meeting Scheduling example in Figure 1 (i*) and Figure 3 (BPMN). A new requirement within in the form of a task
dependency between the Meeting Initiator (i.e. the dependee) and the Meeting Scheduler (i.e. the depender) to ProvideParticipantPrioritization. Participant prioritization
means that the Meeting Initiator must now prioritize the current list of participants in
order for the Meeting Scheduler to MergeAvailableDates and FindAnAgreeableSlot
effectively.
Given the application of our approach for guiding an analysts decision, it can be
inferred that the effect for ParticipantPrioritizationProvided will propagate within
the i* model as a fulfillment condition on the SchedulerSchedulesMeetingTask. Furthermore, given Consistency Rule 3, requires that ParticipantPrioritizationProvided
occurs prior to the fulfillment of the MergeAvailableDates fulfillment conditions.
This information can then be used to highlight the scope of change within the BPMN
model to a point within a trajectory prior to the required effects of MergeAvailableDates, where an activity controlled by the initiator is able to realize the required effect.
3.5 Constrained Development of a High-level Conceptual Model given a
Business Process Model
The following steps provide systematic guidance for developing an i* model given an
already existing process model. Figure 5, illustrates the constrained development of
the Patient Treatment BPMN model in Figure 2.
− Step 1: Map elements to equivalent constructs within the i* model.
− Step 1.1: Map Participants. Both pools and lanes in a BPMN model represent
actors in an i* model. These can be directly translated into the model.
− Step 1.2: Map Activities. Represent activities and sub-processes as ‘primitively
workable’ tasks assigned to actors in i*.
− Step 2. Apply intentional reasoning.
− Step 2.1: Query the Intention of Tasks. Intentional reasoning is applied to identify higher-level intentional elements and dependencies by querying the intention of tasks. This step aims to guide the further understanding and representation of an actors motivations.
− Step 2.2: Query the Intention of Flow-Links. Analyze control and message
flow between actor boundaries to identify goal, task and resource dependencies.
These types of links can be used as a primary heuristic for identifying possible
dependencies between actors.
− Step 3: Identify soft-goal dependencies in the i* model. The representation of softgoals (including dependencies) are not in the scope of the BPMN notation.
3.5.1 Reflecting Changes in a BPMN Model to an associated i* Model
These steps indicate how BPMN model change may be reflected in the i* model:
− Step 1: For each classification of change, apply the following changes.
− Addition of a swimlane or pool. If a swimlane or pool is added, then a new actor will be required within the i* model. This will include the addition of new
dependencies and tasks within the i* model. A primary heuristic for identifying
dependencies includes message flow links and control flow links between pools
and lanes (message flow ndicates a resource dependency for some information).
− Addition a task to an existing swimlane or pool. If a new task is added to a
swimlane or pool, this will require a task to be decomposed from the root node
of the routine traceable to the current process.
− Step 2: Re-apply consistency rules assess whether maintenance.
Fig. 4. An i* ‘Patient Treatment’ Process
Consider now a scenario where the business process model is modified to improve
the performance of the IssuePrescription task which has been identified to be a major
operational bottleneck. The task is improved by including a task before hand which
checks the patient’s previous medical history to identify previous prescriptions for the
patient for similar illnesses (e.g. common flu). We name the task CheckPatientMedicalHistory. Furthermore, the client is now encouraged to provide information on his
medical background, which we represent as a task named ProvideMedicalHistoryInformation. We now proceed to add an additional task within the bounds of the Doctor
agent and an additional task within the bounds of the Patient agent.
Step 1:
Pools and
Lanes as Actors;
Activities &
Sub-Processes
Step 2:
Querying
Intentions –
Tasks and
Flow Links
Step 3:
Soft-Goals
As in the previous case we use intentional reasoning to identify that the added task,
within the Doctor agent, contributes to the higher level task of TreatingPatients. We
apply the same technique to justify the placement of the ProvideMedicalHistoryInformation task as a decomposition task under the RequestMedicine task.
The added message flow in the BPMN diagram is represented as a resource dependency between the Patient and the Doctor, where the Doctor requires the Patient
to provide his previous medical history. We also introduce the soft-goal between the
Patient and the Doctor, titled TimelyDrugPrescription, indicating the fact that the
Doctor will try to improve the time required to prescribe medication to the Patient.
4 Conclusion
In this work, we have illustrated an initial approach for supporting the lifecycle of
business process models with the complementary use of i* - a well developed notation for modeling organizational contexts, and BPMN – a newly developed notation
for modeling business processes. The approach for reflecting changes in organizational context to changes in the design of business processes provides an effective
mechanism for aligning business processes with organizational objectives. Similarly,
operational improvements can be mapped back to organizational objectives to facilitate analysis and ensure no conflicts exist with existing objectives. Although these
steps are preliminary we believe their systematic nature makes them available for
automation in all phases, and are pursuing this task, through the development of a
software tool, along with further refinement of the approach.
References
1. Smith, H. Fingar, P.: Business Process Management – The Third Wave. Meghan-Kiffer
Press, Tampa Florida (2003)
2. Hammer, M. Champy, J.: Reengineering the Corporation: A Manifesto for Business Revolution. HarperBusiness, (1993)
3. Dumas, M. Aalst, W. M. P. and Hofstede, A. H.: Process-Aware Information Systems:
Bridging People and Software Through Process Technology. Wiley-Interscience (2005)
4. Bider, I. Johannesson, P.: Tutorial on: Modeling Dynamics of Business Processes – Key for
Building Next Generation of Business Information Systems. In: The 21st International Conference on Conceptual Modeling (ER2002), Tampere, FL, October 7-11 (2002)
5. Kavakli, V. and Loucopoulos, P.: Goal-Driven Business Process Analysis - Application in
Electricity Deregulation. In: Information Systems, 24(3) (1999) 187-207
6. Loucopoulos, P. and Kavakli, E.: Enterprise Modeling and the Teleological Approach to
Requirements Engineering. In: International Journal of Intelligent and Cooperative Information Systems 4(1) (1995) 45-79
7. Katzenstein G. Lerch F. J.: Beneath the surface of organizational processes: a social representation framework for business process redesign. In: ACM Transactions on Information
Systems (TOIS), 18(4) (2000) 383-422
8. Yu, E.: Models for Supporting the Redesign of Organizational Work. In: Proceedings, Conf.
on Organizational Computing Systems (COOCS'95) Milpitas, California, USA, August 13-
16 (1995) 225-236
9. Yu, E.: Modelling Strategic Relationships for Process Reengineering. PhD Thesis, Graduate
Department of Computer Science, University of Toronto, Toronto, Canada (1995) ~124
10.Davies, I, Green, P. Rosemann, M. Gallo, S.: Conceptual Modelling – What and Why in
Current Practice. In: Lecture Notes in Computer Science, Volume 3288 (2004) 30-42
11.Cysneiros, L.M. Yu, E. Addressing Agent Autonomy in Business Process Management -
with Case Studies on the Patient Discharge Process. In: Proc. of the 2004 Information Resources Management Association Conference, New Orleans, May, (2004)
12.Andersson, B. Bergholtz, M. Edirisuriya, A. Ilayperuma, T. Johannesson, P.: A Declarative
Foundation of Process Models. In: Lecture Notes in Computer Science, Volume 3520
(2005) 233–247
13.Krishna, A., Ghose, A. K., Vranesevic, A.: Loosely-coupled Consistency between AgentOriented Conceptual Models and UML Sequence Diagrams. In: Special issue on AgentOriented Software Development Methodologies International Journal of Multi-Agent and
Grid Systems. Accepted. (2006)
14.Dasgupta, A., Salim, F., Krishna, A., Ghose, A. K.: Hybrid Modelling using i* and AgentSpeak (L) Agents in Agent-Oriented Software Engineering. To appear in: The Proceedings of
8th International Conference on Enterprise Information System (ICEIS-2006), Paphos, Cyprus, May 23-27 (2006)
15.Krishna, A., Guan, Y., Sambattheera, C., Ghose, A. K.: Agent-based Prototyping of Webbased Systems. To appear in: The Proceedings of the 19th International Conference on Industrial and Engineering Applications of Artificial Intelligence and Expert Systems (IEAAIE-2006), Springer-Verlag Lecture Notes in Computer Science, Annecy, France, 27-30
June (2006)
16.Koliadis, G. Vranesevic, A. Bhuiyan, M. Krishna, A. Ghose, A.: A Combined Approach for
Supporting the Business Process Model Lifecycle. To appear in: The Proceedings of the
10th Pacific Asia Conference on Information Systems (PACIS), July 6-9, Kuala Lumpur,
Malaysia (2006).
17.White, S. Business Process Modeling Notation (BPMN) Version 1.0. Business Process
Management Initiative (BPMI.org), May (2004)
18.Giorgini, P. Kolp, M. Mylopoulos, J. Pistore, M.: The Tropos Methodology: an overview.
In: Methodologies And Software Engineering For Agent Systems. Kluwer Academic Publishing (2004)
19.Fuxman, A. Liu, L. Mylopoulos, J. Pistore, M. Roveri, M. Traverso, P.: Specifying and
analyzing early requirements in Tropos. In: Requirements Engineering, Springer London,
9(2) (2004) 132–150
20.Fischer, L.: Workflow Handbook 2005. Workflow Management Coalition, (WfMC) (2005)
21.Ouyang, C. W.M.P. van der Aalst, Dumas, M. and ter Hofstede, A.H.M.: Translating
BPMN to BPEL. BPM Center Report BPM-06-02, BPMcenter.org, (2006)
22.Becker, J. Indulska, M. and Rosemann, M. Green, P.: Do Process Modelling Techniques
Get Better? A Comparative Ontological Analysis of BPMN. In: Campbell, Bruce and Underwood, Jim and Bunker, Deborah, Eds. Proceedings 16th Australasian Conference on
Information Systems, Sydney, Australia (2005)
23.Lamsweerde, A. Goal-Oriented Requirements Engineering: A Guided Tour. In: The 5th
International Symp. In Requirements Engineering (RE’01), Aug. (2001)